% =============================================================================
% Section 11 — Persistence & Crash Recovery
% =============================================================================
\section{Persistence and Crash Recovery}
\label{sec:persistence}

Every \pebble{} node persists its state to an isolated data directory
(\texttt{dataDir}). The persistence layer ensures that no acknowledged
mutation is lost across node restarts and that the node can reconstruct its
full in-memory state on startup.

\subsection{On-Disk Layout}

\begin{lstlisting}[caption={Data directory layout.}, language={}]
dataDir/
  identity.json          # Node keypair and nodeId (generated once)
  metadata.json          # Atomic sequence and vector counter
  mutations.log          # Append-only NDJSON mutation log (with fsync)
  snapshots/
    2026-08-20T00-00-00-000Z.snapshot.json
    2026-08-21T00-00-00-000Z.snapshot.json
    ...
\end{lstlisting}

\subsection{Append-Only Mutation Log}

The \texttt{MutationLog} (\texttt{packages/core/src/persistence/log.ts})
appends each accepted local mutation to \texttt{mutations.log} as a single
newline-delimited JSON record (NDJSON format). Each append is followed by an
\texttt{fsync} call to flush the kernel buffer to durable storage before the
mutation is acknowledged to the caller.

\begin{definition}[Durability guarantee]
A mutation $m$ is \emph{durable} at node $n$ if $m$ has been serialised to
\texttt{mutations.log} and the \texttt{fsync} call has returned. Durable
mutations survive process crashes, OS crashes, and power failures (assuming the
storage medium respects \texttt{fsync} semantics).
\end{definition}

\begin{remark}[Gossip-received mutations]
Mutations received via gossip from peer nodes are applied to the in-memory store
and emitted as \texttt{mutation:applied} events. The orchestrator (\texttt{Pebble}
class) listens for these events and appends received mutations to the log as
well, ensuring that any mutation in the in-memory state is also represented
on disk.
\end{remark}

\subsection{Snapshots}

Taking a snapshot serialises the entire current state (including tombstones)
and the current version vector to a JSON file in \texttt{snapshots/}:

\begin{definition}[Snapshot]
A \emph{snapshot} $\mathcal{S}$ is a tuple
\[
  \mathcal{S} = \langle \mathit{version},\; t_\mathcal{S},\; V,\;
                        \mathit{nextSeq},\; S \rangle
\]
where $V$ is the version vector at snapshot time, $\mathit{nextSeq}$ is the
next local sequence number, and $S$ is the current state map
(key $\to$ winning mutation, including tombstones).
\end{definition}

Snapshot filenames are ISO 8601 timestamps safe for all operating systems
(colons replaced with hyphens). Lexicographic sort of filenames corresponds to
chronological order, enabling simple latest-snapshot selection.

Snapshots serve two purposes:
\begin{enumerate}[leftmargin=1.5em]
  \item \textbf{Faster recovery}: recovery needs only to replay log entries
        after the snapshot timestamp, rather than the full history.
  \item \textbf{Log truncation (future)}: once a snapshot is taken, older log
        entries that predate the snapshot can safely be discarded (\cref{sec:future}).
\end{enumerate}

\subsection{Crash Recovery Pipeline}

On node startup, the \texttt{recover} function reconstructs in-memory state
from the on-disk data:

\begin{algorithm}[H]
\caption{Node crash recovery (\texttt{recover})}
\label{alg:recovery}
\begin{algorithmic}[1]
\Require Node \texttt{dataDir}, empty \texttt{store}, \texttt{nodeId}
\State $\mathcal{S} \leftarrow \texttt{loadLatestSnapshot(dataDir)}$
\State $\mathcal{M} \leftarrow \texttt{loadMetadata(dataDir)}$
\State $\mathit{nextSeq} \leftarrow \mathcal{M}.\mathit{nextLocalSequence}$

\If{$\mathcal{S} \neq \bot$}
  \State $\texttt{store.importSnapshot}(\mathcal{S}.S,\; \mathcal{S}.V,\; \mathcal{S}.\mathit{nextSeq})$
  \State $\mathit{nextSeq} \leftarrow \max(\mathit{nextSeq},\; \mathcal{S}.\mathit{nextSeq})$
\Else
  \State $\texttt{store.setNextLocalSequence}(\mathit{nextSeq})$
\EndIf

\State $\mathit{log} \leftarrow \texttt{MutationLog(dataDir).readAll()}$
\For{each mutation $m$ in $\mathit{log}$}
  \State $\texttt{store.applyMutation}(m,\; \{\mathit{skipSignatureVerification}: \top\})$
  \If{$m.\mnid = \mathit{nodeId}$}
    \State $\mathit{nextSeq} \leftarrow \max(\mathit{nextSeq},\; m.\mseq + 1)$
  \EndIf
\EndFor

\State $\texttt{store.setNextLocalSequence}(\mathit{nextSeq})$
\end{algorithmic}
\end{algorithm}

\paragraph{Correctness argument.}
The log replay is idempotent: duplicate mutations (e.g., those already covered
by the snapshot) are silently skipped by the store's duplicate detection.
Mutations that arrived after the snapshot timestamp but before the crash are
replayed in full. The $\mathit{nextSeq}$ watermark is advanced to be strictly
greater than any local sequence seen, preventing sequence number reuse across
restarts (Invariant~\ref{inv:mono-seq}).

\subsection{Corrupted Log Tail}

If the node crashed mid-write, the last line of \texttt{mutations.log} may be
a truncated, invalid JSON record. The log reader handles this by discarding any
trailing partial line (lines that do not parse as valid mutations are skipped
with a warning). Since the line was not yet synced (the process crashed before
\texttt{fsync} returned), it was never acknowledged, and discarding it
preserves the durability guarantee.
